% GENERATED FILE. Edit canonical lessons, then run npm run generate:books.
% canonical-source-hash: fnv1a64:406304b680622841
% canonical-lessons: MW-C205-tikat, MW-C205-chukno, MW-C205-bharo, MW-C205-atakno, MW-C205-kandaktar

\chapter{Missing the Bus}
\label{ch:mw-missing-the-bus}
\hlchaptermodality{205}

\begin{chapteropening}
\textbf{By the end of this chapter:} \emph{I can ask the fare and say I missed the bus or got stuck.}

Complete the last lesson of chapter 205: I can ask the fare and say I missed the bus or got stuck.
\end{chapteropening}

\section[\texorpdfstring{\mw{टिकट} (ṭikaṭ)}{टिकट (ṭikaṭ)}]{\mw{टिकट} (ṭikaṭ) — a ticket}

\label{lesson:MW-C205-tikat}



\begin{quote}
Before the new one: say the Marwadi for an edge, then the Marwadi for a map.
\end{quote}

\subsection*{You'll want to know: \mw{टिकट}}
\textbf{\mw{टिकट}} — \emph{ṭikaṭ} — \textquotedblleft{}a ticket\textquotedblright{}.

\begin{quote}
\mw{टिकट} \mw{कठै} \mw{मिलै} \mw{है}?
\end{quote}

\begin{quote}
\emph{ṭikaṭ kaṭhai milai hai?}
\end{quote}

\begin{quote}
Where can I get a ticket?
\end{quote}

A masculine noun.

\begin{cousinweb}
\textbf{\mw{टिकट}} (\emph{ṭikaṭ}) is borrowed from English \emph{ticket}.
\end{cousinweb}

\begin{grammarlens}[title={What you've built}]
One more word for this chapter.
\end{grammarlens}

\subsection*{Your turn}
\begin{itemize}
\raggedright
  \item \emph{Say it:} \emph{ṭikaṭ}
  \item \emph{Say it:} \emph{ṭikaṭ}, once more
  \item \emph{Say it:} \emph{ṭikaṭ kaṭhai milai hai?} — Where can I get a ticket?
  \item \emph{Recall:} say the Marwadi for an edge, then the Marwadi for a map, then say \emph{ṭikaṭ} again
\end{itemize}

\begin{tcolorbox}[breakable,colback=teal!4,colframe=teal!35!black,title={Before you move on}]
What does \mw{टिकट} mean? (\textquotedblleft{}A ticket\textquotedblright{}.) Say it once more.
\end{tcolorbox}
\section[\texorpdfstring{\mw{चूकणो} (chūkṇo)}{चूकणो (chūkṇo)}]{\mw{चूकणो} (chūkṇo) — to miss (a bus); to slip up}

\label{lesson:MW-C205-chukno}



\begin{quote}
Before the new one: say the Marwadi for a map, then the Marwadi for a ticket.
\end{quote}

\subsection*{You'll want to know: \mw{चूकणो}}
\textbf{\mw{चूकणो}} — \emph{chūkṇo} — \textquotedblleft{}to miss (a bus); to slip up\textquotedblright{}.

\begin{quote}
\mw{हूं} \mw{रोज} \mw{बस} \mw{चूकूं} \mw{हूं।}
\end{quote}

\begin{quote}
\emph{hū̃ roj bas chūkū̃ hū̃.}
\end{quote}

\begin{quote}
I miss the bus every day.
\end{quote}

Past (I …): \mw{हूं} \mw{चूक्यो}, \emph{hū̃ chūkyo}.

\begin{cousinweb}
\textbf{\mw{चूकणो}} (\emph{chūkṇo}) is inherited from a reconstructed \emph{*cukk}, \textquotedblleft{}fall short of, stop\textquotedblright{}; CDIAL lists Marwari \emph{cūkṇo}, \textquotedblleft{}to be completed\textquotedblright{}, and Gujarati \emph{cukvũ}, \textquotedblleft{}to fail, miss\textquotedblright{} (CDIAL 4848). \textbf{\mw{चुकावणो}}, \textquotedblleft{}to pay\textquotedblright{}, is its causative.
\end{cousinweb}

\begin{grammarlens}[title={What you've built}]
One more word for this chapter.
\end{grammarlens}

\subsection*{Your turn}
\begin{itemize}
\raggedright
  \item \emph{Say it:} \emph{chūkṇo}
  \item \emph{Say it:} \emph{chūkṇo}, once more
  \item \emph{Say it:} \emph{hū̃ roj bas chūkū̃ hū̃.} — I miss the bus every day.
  \item \emph{Recall:} say the Marwadi for a map, then the Marwadi for a ticket, then say \emph{chūkṇo} again
\end{itemize}

\begin{tcolorbox}[breakable,colback=teal!4,colframe=teal!35!black,title={Before you move on}]
What does \mw{चूकणो} mean? (\textquotedblleft{}To miss (a bus); to slip up\textquotedblright{}.) Say it once more.
\end{tcolorbox}
\section[\texorpdfstring{\mw{भाड़ो} (bhāṛo)}{भाड़ो (bhāṛo)}]{\mw{भाड़ो} (bhāṛo) — a fare; hire}

\label{lesson:MW-C205-bharo}



\begin{quote}
Before the new one: say the Marwadi for a ticket, then the Marwadi for to miss.
\end{quote}

\subsection*{You'll want to know: \mw{भाड़ो}}
\textbf{\mw{भाड़ो}} — \emph{bhāṛo} — \textquotedblleft{}a fare; hire\textquotedblright{}.

\begin{quote}
\mw{भाड़ो} \mw{कितणो} \mw{है}?
\end{quote}

\begin{quote}
\emph{bhāṛo kitṇo hai?}
\end{quote}

\begin{quote}
How much is the fare?
\end{quote}

A masculine noun.

\begin{cousinweb}
\textbf{\mw{भाड़ो}} (\emph{bhāṛo}) is inherited from a reconstructed \emph{*bhārta}, \textquotedblleft{}pertaining to servants\textquotedblright{}, by way of Middle Indo-Aryan \emph{bhāṭa}, \textquotedblleft{}wages\textquotedblright{}. CDIAL gives no Marwari form in this entry, but lists Sindhi \emph{bhāṛo}, \textquotedblleft{}hire\textquotedblright{}, and Hindi \emph{bhāṛā}, \textquotedblleft{}hire, fare, rent\textquotedblright{} (CDIAL 9468).
\end{cousinweb}

\begin{grammarlens}[title={What you've built}]
One more word for this chapter.
\end{grammarlens}

\subsection*{Your turn}
\begin{itemize}
\raggedright
  \item \emph{Say it:} \emph{bhāṛo}
  \item \emph{Say it:} \emph{bhāṛo}, once more
  \item \emph{Say it:} \emph{bhāṛo kitṇo hai?} — How much is the fare?
  \item \emph{Recall:} say the Marwadi for a ticket, then the Marwadi for to miss, then say \emph{bhāṛo} again
\end{itemize}

\begin{tcolorbox}[breakable,colback=teal!4,colframe=teal!35!black,title={Before you move on}]
What does \mw{भाड़ो} mean? (\textquotedblleft{}A fare; hire\textquotedblright{}.) Say it once more.
\end{tcolorbox}
\section[\texorpdfstring{\mw{अटकणो} (aṭakṇo)}{अटकणो (aṭakṇo)}]{\mw{अटकणो} (aṭakṇo) — to get stuck}

\label{lesson:MW-C205-atakno}



\begin{quote}
Before the new one: say the Marwadi for to miss, then the Marwadi for a fare.
\end{quote}

\subsection*{You'll want to know: \mw{अटकणो}}
\textbf{\mw{अटकणो}} — \emph{aṭakṇo} — \textquotedblleft{}to get stuck\textquotedblright{}.

\begin{quote}
\mw{बस} \mw{भीड़} \mw{मांय} \mw{अटकै} \mw{है।}
\end{quote}

\begin{quote}
\emph{bas bhīṛ mā̃y aṭkai hai.}
\end{quote}

\begin{quote}
The bus gets stuck in the crowd.
\end{quote}

Past (I …): \mw{हूं} \mw{अटक्यो}, \emph{hū̃ aṭakyo}.

\begin{cousinweb}
\textbf{\mw{अटकणो}} (\emph{aṭakṇo}) is inherited from a reconstructed \emph{*aṭṭakk}, \textquotedblleft{}obstruct, stop\textquotedblright{}; CDIAL lists Marwari \emph{aṭakṇo}, \textquotedblleft{}to remain\textquotedblright{}, and Hindi \emph{aṭaknā}, \textquotedblleft{}to stick\textquotedblright{} (CDIAL 182).
\end{cousinweb}

\begin{grammarlens}[title={What you've built}]
One more word for this chapter.
\end{grammarlens}

\subsection*{Your turn}
\begin{itemize}
\raggedright
  \item \emph{Say it:} \emph{aṭakṇo}
  \item \emph{Say it:} \emph{aṭakṇo}, once more
  \item \emph{Say it:} \emph{bas bhīṛ mā̃y aṭkai hai.} — The bus gets stuck in the crowd.
  \item \emph{Recall:} say the Marwadi for to miss, then the Marwadi for a fare, then say \emph{aṭakṇo} again
\end{itemize}

\begin{tcolorbox}[breakable,colback=teal!4,colframe=teal!35!black,title={Before you move on}]
What does \mw{अटकणो} mean? (\textquotedblleft{}To get stuck\textquotedblright{}.) Say it once more.
\end{tcolorbox}
\section[\texorpdfstring{\mw{कंडक्टर} (kaṇḍakṭar)}{कंडक्टर (kaṇḍakṭar)}]{\mw{कंडक्टर} (kaṇḍakṭar) — a bus conductor}

\label{lesson:MW-C205-kandaktar}



\begin{quote}
Before the new one: say the Marwadi for a fare, then the Marwadi for to get stuck.
\end{quote}

\subsection*{You'll want to know: \mw{कंडक्टर}}
\textbf{\mw{कंडक्टर}} — \emph{kaṇḍakṭar} — \textquotedblleft{}a bus conductor\textquotedblright{}.

\begin{quote}
\mw{कंडक्टर} \mw{टिकट} \mw{देवै} \mw{है।}
\end{quote}

\begin{quote}
\emph{kaṇḍakṭar ṭikaṭ devai hai.}
\end{quote}

\begin{quote}
The conductor hands out tickets.
\end{quote}

A masculine or feminine noun, by the person meant.

\begin{cousinweb}
\textbf{\mw{कंडक्टर}} (\emph{kaṇḍakṭar}) is borrowed from English \emph{conductor}.
\end{cousinweb}

\begin{grammarlens}[title={What you've built}]
One more word for this chapter.
\end{grammarlens}

\subsection*{Your turn}
\begin{itemize}
\raggedright
  \item \emph{Say it:} \emph{kaṇḍakṭar}
  \item \emph{Say it:} \emph{kaṇḍakṭar}, once more
  \item \emph{Say it:} \emph{kaṇḍakṭar ṭikaṭ devai hai.} — The conductor hands out tickets.
  \item \emph{Recall:} say the Marwadi for a fare, then the Marwadi for to get stuck, then say \emph{kaṇḍakṭar} again
\end{itemize}

\begin{tcolorbox}[breakable,colback=teal!4,colframe=teal!35!black,title={Before you move on}]
What does \mw{कंडक्टर} mean? (\textquotedblleft{}A bus conductor\textquotedblright{}.) Say it once more.
\end{tcolorbox}
